% Figure: grid frequency response to a sudden generator trip. Frequency
% (Hz) vs time (s) showing the nadir, primary response arrest, and
% secondary recovery to nominal. Nominal 50 Hz.
\begin{figure}[htbp]
\centering
\begin{tikzpicture}
\begin{axis}[
  width=13cm, height=8cm,
  xlabel={time after trip (s)},
  ylabel={grid frequency (Hz)},
  xmin=0, xmax=60, ymin=49.5, ymax=50.1,
  xtick={0,10,20,30,40,50,60},
  ytick={49.5,49.6,49.7,49.8,49.9,50.0,50.1},
  axis lines=left,
  tick label style={font=\scriptsize},
  label style={font=\small},
]
% nominal line
\draw[enggray, dashed] (axis cs:0,50) -- (axis cs:60,50);
\node[enggray, font=\scriptsize, anchor=south east] at (axis cs:60,50) {nominal 50 Hz};
% the trip happens at t=2; frequency falls to nadir, arrested, recovers
\addplot[engblue, very thick, smooth] coordinates {
  (0,50.0)(2,50.0)(3,49.85)(4,49.72)(5,49.66)(6,49.65)(8,49.68)(12,49.74)(20,49.82)(30,49.90)(45,49.97)(60,50.0)
};
% nadir marker
\filldraw[engred] (axis cs:6,49.65) circle [radius=1.6pt];
\node[engred, font=\scriptsize, anchor=north] at (axis cs:6.5,49.64) {nadir};
% phase annotations
\draw[engamber, |-|] (axis cs:2,49.55) -- (axis cs:8,49.55);
\node[engamber, font=\scriptsize, anchor=north] at (axis cs:5,49.55) {primary (inertia + governor)};
\draw[englime!60!black, |-|] (axis cs:8,49.55) -- (axis cs:60,49.55);
\node[englime!60!black, font=\scriptsize, anchor=north] at (axis cs:34,49.55) {secondary (AGC)};
% trip arrow
\draw[engred, -{Stealth}] (axis cs:2,49.95) -- (axis cs:2,50.0);
\node[engred, font=\scriptsize, anchor=north] at (axis cs:2,49.93) {trip};
\end{axis}
\end{tikzpicture}
\caption[Grid frequency response to a generator trip]{Grid frequency
after a sudden generator trip. The rotor inertia of the synchronous fleet
limits the initial rate of fall; governor (primary) response arrests the
decline at the nadir within seconds; automatic generation control
(secondary) restores the frequency to nominal over tens of seconds to
minutes. A shallower nadir means more inertia and faster primary
response. Illustrative trajectory.}
\label{fig:vol04ch14:frequency-response}
\end{figure}
